% Originally: content/week-13.tex, \section{Stokes' theorem} (Corsin Week 13)


\section{Stokes' theorem}
\label{sec:stokes_general}

% Source: Corsin Nick/Class Notes/Week 13.pdf, p. 2
\begin{theorem}[Stokes' theorem]
\label{thm:stokes_general}
Let $M \subseteq \mathbb{R}^n$ be an orientable $m$-manifold with (possibly empty) boundary,
and let $\omega \in \Omega^{m-1}(\mathbb{R}^n)$ have compact support in $\mathbb{R}^n$. Then
\[\int_{\partial M} \omega = \int_M d\omega.\]
\end{theorem}

\begin{ainote}
\Cref{thm:stokes_classical} was stated ahead of Corsin's own notes, as background
needed to solve \cref{ex:12.3}, in its classical $\mathbb{R}^3$ curl-theorem form. It is
exactly the case $m=2$, $n=3$, $\omega=\omega_F$ of the general theorem above, as
\cref{sec:stokes_in_r3} below confirms directly from Corsin's own derivation.
\end{ainote}

\subsection{Example: Green's theorem}
\label{sec:green_theorem_example}

% Source: Corsin Nick/Class Notes/Week 13.pdf, p. 2
\begin{example}[The exterior derivative of a $1$-form on $\mathbb{R}^2$]
\label{ex:exterior_derivative_r2}
Let $\omega(x,y) = P(x,y)\,dx + Q(x,y)\,dy$ be a $1$-form on $\mathbb{R}^2$. Then
\[
d\omega = dP\wedge dx + dQ\wedge dy = \left(\frac{\partial P}{\partial x}dx +
  \frac{\partial P}{\partial y}dy\right)\wedge dx + \left(\frac{\partial Q}{\partial x}dx +
  \frac{\partial Q}{\partial y}dy\right)\wedge dy = \left(\frac{\partial Q}{\partial x} -
  \frac{\partial P}{\partial y}\right) dx\wedge dy,
\]
using $dx\wedge dx = dy\wedge dy = 0$ and $dy\wedge dx = -dx\wedge dy$.
\end{example}

\subsection{Example: proof of Green's theorem}
\label{sec:green_theorem_proof}

% Source: Corsin Nick/Class Notes/Week 13.pdf, p. 3
\begin{ainote}
The theorem below is reproduced by Corsin directly from the lecture script (numbered there as
Theorem 14.24), as with the two definitions quoted in \cref{sec:orientable_manifolds}.
\end{ainote}

\begin{theorem}[Green's theorem in the plane]
\label{thm:green_theorem}
Let $F : U \to \mathbb{R}^2$ be a continuously differentiable vector field on an open set
$U \subseteq \mathbb{R}^2$. Then, for any compact $C^1$ domain $\Omega \subseteq U$,
\[\int_\Omega \curl F\,dx = \int_{\partial\Omega} F\cdot\tau\,dL,\]
with $\tau = \nu^\perp = (-\nu_2,\nu_1)$, where $\nu$ is the exterior unit normal.
\end{theorem}

Writing $F(x,y) = (P(x,y),Q(x,y))$, so that $\curl F = \partial_xQ - \partial_yP$ (the scalar
$2$-dimensional curl), \cref{thm:green_theorem} follows from
\cref{thm:stokes_general,ex:exterior_derivative_r2} applied to $\omega = P\,dx+Q\,dy$:
\[
\int_\Omega \curl F\,dx\,dy = \int_\Omega \left(\frac{\partial Q}{\partial x} -
  \frac{\partial P}{\partial y}\right) dx\wedge dy = \int_\Omega d\omega
  \overset{\text{Stokes}}{=} \int_{\partial\Omega} P(x,y)\,dx + Q(x,y)\,dy
  = \int_{\partial\Omega} F\cdot\tau\,dL.
\]

\begin{proof}[Justification of the last step]
Suppose $\gamma : [0,L] \to \partial\Omega$ is a $C^1$ curve of unit speed and positive
orientation, i.e.\ $\gamma'(t) = \tau(\gamma(t))$. Then
\[\langle F,\tau\rangle = \langle F,\gamma'(t)\rangle = (P\,dx+Q\,dy)(\gamma'(t)),\]
since $\tau = \begin{pmatrix}0&-1\\1&0\end{pmatrix}\nu$ is the tangent vector consistent with
the induced orientation, so integrating $(P\,dx+Q\,dy)(\gamma'(t))$ over $t \in [0,L]$ is
exactly $\int_{\partial\Omega} P\,dx+Q\,dy$ by \cref{def:integral_of_form_over_submanifold}.
\end{proof}

\subsection{Example: Stokes' theorem in \texorpdfstring{$\mathbb{R}^3$}{R\^{}3}}
\label{sec:stokes_in_r3}

% Source: Corsin Nick/Class Notes/Week 13.pdf, p. 4
Recall from \cref{ex:exterior_derivative_gradient_curl} that for $\omega_F =
\langle F,\cdot\rangle$, $d\omega_F = \langle\curl F,\cdot\times\cdot\rangle$. So, if $\Omega
\subseteq \mathbb{R}^3$ is a $2$-dimensional submanifold with boundary, parametrized by $\phi$
with tangent vectors $\tau_i := \partial_i\phi$ and exterior unit normal $\nu$, then
\[
\int_\Omega d\omega_F = \int_\Omega \langle\curl F,\tau_1\times\tau_2\rangle\,dx
  = \int_\Omega \langle\curl F,\nu\rangle\,d\vol_2 = \int_\Omega \vec\nabla\times\vec F\cdot
  d\vec A,
\]
and
\[\int_{\partial\Omega} \omega_F = \int_{\partial\Omega} \langle F,\tau\rangle\,d\vol_1
  = \int_{\partial\Omega} \vec F\cdot d\vec s.\]
By \cref{thm:stokes_general}, these are equal, which recovers \cref{thm:stokes_classical}.

% Source: Jérôme Paschoud/Notizen/Woche 13.1 Stokes.pdf, p. 4
\begin{theorem}[Stokes' theorem in 3D]
\label{thm:stokes_3d}
For an orientable, parametrized $2$-dimensional submanifold $M \subseteq \mathbb{R}^3$ with boundary $\partial M$, and a $C^1$ vector field $F$:
\[\int_M \curl(F) \cdot \nu \,d\vol_2 = \int_{\partial M} F \cdot \tau \,d\vol_1.\]
The orientation of $\partial M$ is chosen such that when traversing $\partial M$ in the direction of $\tau$, the surface $M$ (given by the normal $\nu$) is to the left (the ``right-hand rule'').
\end{theorem}

% Generator: Claude Opus 5 (High) --- FIG-26-03. This section is the chapter's main one and had
% no figure at all, while the theorem just above states its orientation convention purely in
% words. "The surface is to the left" is the one convention in this part of the course that a
% reader cannot reliably reconstruct from prose, and getting it backwards flips the sign of
% every answer, which is why it earns a picture here rather than a longer sentence.
%
% GEOMETRY, checked. The surface is a horizontal disc in the plane z = 0, drawn under the
% projection screen = (x, 0.4y + z), so a disc of radius 2 becomes an ellipse with semi-axes
% 2 and 0.8. The parametrization (2 cos t, 2 sin t, 0) maps to screen (2 cos t, 0.8 sin t),
% which runs right -> top -> left -> bottom as t increases, i.e. COUNTERCLOCKWISE on screen,
% and counterclockwise seen from +z is the sense that goes with the upward normal.
%
% The marked point is t = pi/2, the top of the ellipse, chosen because it is the only place
% where all three vectors are visually distinct:
%   tau     = d/dt (2 cos t, 2 sin t, 0) at t = pi/2, normalized  = (-1, 0, 0) -> LEFT on screen
%   nu      = (0, 0, 1)                                                        -> UP on screen
%   inward  = nu x tau = (0,0,1) x (-1,0,0) = (0, -1, 0)                       -> DOWN on screen
% The cross product is the content of the figure and is worth checking rather than trusting:
% (0,0,1) x (-1,0,0) = (0*0 - 1*0, 1*(-1) - 0*0, 0*0 - 0*(-1)) = (0,-1,0), and at the point
% (0,2,0) the direction (0,-1,0) does point back towards the centre, so it really is inward.
% Read that as the walker's rule: facing tau with nu upright, the left hand points along
% nu x tau, which is into M. At t = 0 this figure would be undrawable, since tau and nu would
% both project onto the same upward screen direction.
\begin{center}
\begin{tikzpicture}[>=Latex, scale=1.1]
  % The surface, and its boundary traversed counterclockwise on screen.
  \fill[MidnightBlue!7] (0,0) ellipse (2 and 0.8);
  \draw[MidnightBlue!70, thick] (0,0) ellipse (2 and 0.8);
  \foreach \a in {150,215,330} {
    \draw[MidnightBlue!70, very thick, ->]
      ({2*cos(\a-8)},{0.8*sin(\a-8)}) arc[start angle=\a-8, end angle=\a+8,
        x radius=2, y radius=0.8];
  }
  \node[MidnightBlue!70, font=\small] at (1.15,-0.30) {$M$};
  \node[MidnightBlue!70, font=\scriptsize, anchor=north] at (-1.62,-0.62) {$\partial M$};

  % The three directions at the top of the rim, where all three are distinguishable.
  \draw[OliveGreen, very thick, ->] (0,0.8) -- (0,1.85)
    node[anchor=south, font=\scriptsize] {$\nu$};
  \draw[BrickRed, very thick, ->] (0,0.8) -- (-1.05,0.8)
    node[anchor=south, font=\scriptsize] {$\tau$};
  % Anchored to the LEFT of the dashed arrow: anchored to its right it came to rest just before
  % the "$M$" label inside the ellipse, and the two read as one string, "into M M".
  \draw[gray!50!black, thick, dashed, ->] (0,0.8) -- (0,0.12)
    node[anchor=north east, font=\scriptsize, text=gray!45!black] {into $M$};
  \fill[MidnightBlue!70] (0,0.8) circle (1.5pt);

  \node[font=\scriptsize, anchor=north, align=center] at (0,-1.15)
    {walking along $\tau$ with $\nu$ upright, $M$ lies to the left:\\
     the leftward direction is $\nu\times\tau$};
\end{tikzpicture}
\end{center}

\begin{remark}[The convention as a formula]
\label{rem:right_hand_rule_as_cross_product}
\qt{To the left} is a reliable instruction only once one knows which way is up, and $\nu$ is what
supplies that. The picture above turns the sentence into an identity that can be checked instead
of visualised: at a boundary point, the direction pointing from $\partial M$ into $M$ is
\[\text{inward} = \nu\times\tau.\]
Equivalently $\tau = \text{inward}\times\nu$, so any two of the three determine the third. On the
unit disc in the $xy$-plane with $\nu = e_3$, at the point $(1,0,0)$, the positively oriented
$\tau$ is $(0,1,0)$ and $\nu\times\tau = (-1,0,0)$ does point at the origin, confirming the
familiar counterclockwise convention.

Reversing $\nu$ reverses $\tau$, which is why both sides of \cref{thm:stokes_3d} change sign
together and the theorem is insensitive to the choice. What it is \emph{not} insensitive to is
choosing the two independently, and that is the error the formula above rules out.
\end{remark}

\begin{example}[Stokes' theorem on a hemisphere]
\label{ex:stokes_hemisphere_linear_field}
% Michael's Bsp. 20.1.2
Let $F(x,y,z) := (z, x, y)\transp$. We want to compute the flux of the curl of $F$ through the upper hemisphere $H := \{(x,y,z) \in \mathbb{R}^3 \mid x^2 + y^2 + z^2 = 1,\ z \geq 0\}$.

First, the curl of $F$ is:
\[\curl(F) = \begin{pmatrix} \partial_y F_3 - \partial_z F_2 \\ \partial_z F_1 - \partial_x F_3 \\ \partial_x F_2 - \partial_y F_1 \end{pmatrix} = \begin{pmatrix} 1 - 0 \\ 1 - 0 \\ 1 - 0 \end{pmatrix} = \begin{pmatrix} 1 \\ 1 \\ 1 \end{pmatrix}.\]
By \cref{thm:stokes_3d}, the flux of $\curl(F)$ through $H$ equals the line integral of $F$ along the boundary $\partial H$. The boundary is the unit circle in the $xy$-plane, which we parametrize by $\gamma(t) = (\cos t, \sin t, 0)\transp$ for $t \in [0, 2\pi]$. The tangent vector is $\gamma'(t) = (-\sin t, \cos t, 0)\transp$.
The line integral is:
\begin{align*}
\int_{\partial H} F \cdot d\gamma &= \int_0^{2\pi} F(\gamma(t)) \cdot \gamma'(t) \,dt \\
&= \int_0^{2\pi} \begin{pmatrix} 0 \\ \cos t \\ \sin t \end{pmatrix} \cdot \begin{pmatrix} -\sin t \\ \cos t \\ 0 \end{pmatrix} \,dt \\
&= \int_0^{2\pi} \cos^2 t \,dt = \pi.
\end{align*}
Thus, $\int_H \curl(F) \cdot \nu \,d\vol_2 = \pi$.
\end{example}

% Supplement: Analysis_II_Script_v1.pdf, p. 163
\begin{exercise}[Flux through a cylinder, direct and via Stokes]
\label{ex:cylinder_flux_direct_and_stokes}
Let $F(x,y,z) := (yz,x^2,1)\transp$ on $\mathbb{R}^3$, and let $M$ be the lateral surface of the
cylinder $M := \{(x,y,z) : x^2+y^2=1,\ 0<z<1\}$, with the outward-pointing normal
$\nu(x,y,z)=(x,y,0)$. Compute $\int_M\curl(F)\cdot\nu\,d\vol_2$ directly, and again using
Stokes' theorem.
\end{exercise}

\begin{exercisesolution}[Proof of \cref{ex:cylinder_flux_direct_and_stokes}]
% Generator: Claude Sonnet 5 (High)
First, $\curl(F) = (\partial_yF_3-\partial_zF_2,\ \partial_zF_1-\partial_xF_3,\
\partial_xF_2-\partial_yF_1) = (0,\ y,\ 2x-z)$.

\textbf{Directly.} Parametrize $M$ by $\phi(\theta,z) := (\cos\theta,\sin\theta,z)$,
$(\theta,z)\in[0,2\pi)\times(0,1)$; then
$\partial_\theta\phi\times\partial_z\phi = (\cos\theta,\sin\theta,0) = \nu(\phi(\theta,z))$,
already a unit vector, so $d\vol_2 = d\theta\,dz$. Then
$\curl(F)(\phi(\theta,z)) = (0,\sin\theta,2\cos\theta-z)$, and
\[\int_M\curl(F)\cdot\nu\,d\vol_2 = \int_0^1\!\int_0^{2\pi}
  \big\langle(0,\sin\theta,2\cos\theta-z),(\cos\theta,\sin\theta,0)\big\rangle\,d\theta\,dz
  = \int_0^1\!\int_0^{2\pi}\sin^2\theta\,d\theta\,dz = \pi.\]

\textbf{Via Stokes.} The boundary $\partial M$ consists of the two circles $z=0$ and $z=1$; the
orientation induced by the outward normal $\nu$ (right-hand rule) traverses the \emph{bottom}
circle ($z=0$) with $\theta$ increasing and the \emph{top} circle ($z=1$) with $\theta$
\emph{decreasing}. The parameter rectangle $[0,2\pi]\times[0,1]$ has $z=0$ as part of its
positively-oriented boundary traversed left to right and $z=1$ traversed right to left, and the
two vertical edges $\theta=0,2\pi$ cancel since they are the same curve in $M$. This is the
phenomenon drawn in \cref{sec:submanifolds_with_boundary}: one consistently oriented cylinder
induces senses on its two rims that turn opposite ways when both are viewed from the same side,
and \cref{rem:cylinder_opposite_boundary_senses} explains why that is the convention working
rather than failing. With
$\gamma_0(\theta):=(\cos\theta,\sin\theta,0)$ and $\gamma_1(\theta):=(\cos\theta,\sin\theta,1)$,
both for $\theta\in[0,2\pi]$,
\[F(\gamma_0(\theta)) = (0,\cos^2\theta,1), \qquad F(\gamma_1(\theta)) = (\sin\theta,\cos^2\theta,1),\]
and $\gamma_0'(\theta)=\gamma_1'(\theta)=(-\sin\theta,\cos\theta,0)$, so
\[\int_0^{2\pi} F(\gamma_0(\theta))\cdot\gamma_0'(\theta)\,d\theta = \int_0^{2\pi}\cos^3\theta\,d\theta = 0,\]
\[\int_0^{2\pi} F(\gamma_1(\theta))\cdot\gamma_1'(\theta)\,d\theta
  = \int_0^{2\pi}\big({-}\sin^2\theta+\cos^3\theta\big)\,d\theta = -\pi.\]
The bottom circle is traversed with $\theta$ increasing (contributing $+0$) and the top circle
with $\theta$ \emph{decreasing} (contributing $-({-}\pi) = \pi$), so
$\int_{\partial M}F\cdot\tau\,d\vol_1 = 0+\pi = \pi$, agreeing with the direct computation.
\end{exercisesolution}

% Supplement: Analysis_II_Script_v1.pdf, p. 163
\begin{exercise}[Flux through a spherical cap, via Stokes and via Gauss]
\label{ex:spherical_cap_flux_stokes_gauss}
Let $0<h<1$ and let $S := \{(x,y,z)\in S^2 : z>-h\}$ be the part of the unit sphere above
height $-h$. Consider $F(x,y,z) := (-y,x,0)\transp$. Calculate the flux
$\int_S\curl(F)\cdot\nu\,d\vol_2$ (outward normal) directly, once using Stokes' theorem, and
once using the divergence theorem.
\end{exercise}

\begin{exercisesolution}[Proof of \cref{ex:spherical_cap_flux_stokes_gauss}]
% Generator: Claude Sonnet 5 (High)
Here $\curl(F) = (0,0,2)$, a constant field, and on $S^2$ the outward normal is
$\nu(x,y,z)=(x,y,z)$ itself, so $\langle\curl(F),\nu\rangle = 2z$.

\textbf{Directly.} Using colatitude $\varphi\in[0,\arccos(-h)]$ (so $z=\cos\varphi>-h$) and
$d\vol_2=\sin\varphi\,d\varphi\,d\theta$,
\[\int_S 2z\,d\vol_2 = \int_0^{2\pi}\!\int_0^{\arccos(-h)} 2\cos\varphi\sin\varphi\,d\varphi\,d\theta
  = 2\pi\Big[{-}\tfrac12\cos(2\varphi)\Big]_0^{\arccos(-h)} = 2\pi(1-h^2),\]
using $\cos(2\arccos(-h)) = 2h^2-1$.

\textbf{Via Stokes.} $\partial S$ is the circle $x^2+y^2=1-h^2$, $z=-h$, of radius
$r:=\sqrt{1-h^2}$, oriented (right-hand rule, outward normal) by
$\gamma(\theta) := (r\cos\theta,r\sin\theta,-h)$, $\theta\in[0,2\pi]$. Then
$\gamma'(\theta) = (-r\sin\theta,r\cos\theta,0)$ and $F(\gamma(\theta)) =
(-r\sin\theta,r\cos\theta,0)$, so
\[\int_{\partial S}F\cdot d\gamma = \int_0^{2\pi}\big(r^2\sin^2\theta+r^2\cos^2\theta\big)\,d\theta
  = 2\pi r^2 = 2\pi(1-h^2),\]
agreeing with the direct computation.

\textbf{Via Gauss.} Since $\divg(\curl F)\equiv0$ (\cref{ex:curl_grad_div_curl_zero}), the flux
of $\curl(F)$ through any surface with boundary $\partial S$ is the same. Replace $S$ by the
flat disk $D := \{x^2+y^2\leq1-h^2,\ z=-h\}$, with upward normal $(0,0,1)$: $S\cup(-D)$
(the cap together with the disk, reversed) bounds the solid cap $\Omega$, and the divergence
theorem gives $\int_S\curl(F)\cdot\nu\,d\vol_2 - \int_D\curl(F)\cdot(0,0,1)\,d\vol_2 =
\int_\Omega\divg(\curl F)\,d\vol_3 = 0$. Since $\langle\curl F,(0,0,1)\rangle\equiv2$,
\[\int_S\curl(F)\cdot\nu\,d\vol_2 = \int_D 2\,d\vol_2 = 2\cdot\vol_2(D) = 2\pi(1-h^2),\]
the area of a disk of radius $\sqrt{1-h^2}$ times $2$, again the same answer.
\end{exercisesolution}

% Generator: Claude Sonnet 5 (High) --- new content, tying together four theorems this document
% proves separately (FTC-via-Gauss in \cref{sec:gauss_intuition}, Green in
% \cref{sec:green_theorem_proof}, classical Stokes in \cref{sec:stokes_in_r3}, and the
% divergence theorem in \cref{sec:flux_divergence_theorem}) as one theorem, stated once, here.
\begin{remark}[One theorem, four costumes]
\label{rem:one_theorem_four_costumes}
Every integral theorem proved in this part of the course is the single identity
$\int_M d\omega = \int_{\partial M}\omega$ of \cref{thm:stokes_general}, for a different choice
of the dimension $m$ of $M$ and the degree of $\omega$. Collected in one place:
\begin{center}
\small
\begin{tabular}{@{}lllll@{}}
\toprule
\textbf{Theorem} & $M$ & $m$ & $\omega$ & $d\omega$ \\
\midrule
FTC &
  $[a,b]\subset\mathbb{R}$ & $1$ & $f$ ($0$-form) & $f'\,dx$ \\
Green (\cref{thm:green_theorem}) &
  $\Omega\subseteq\mathbb{R}^2$ & $2$ & $P\,dx+Q\,dy$ & $(\partial_xQ{-}\partial_yP)\,dx\!\wedge\! dy$ \\
Classical Stokes (\cref{thm:stokes_classical}) &
  surface $S\subset\mathbb{R}^3$ & $2$ & $\omega_F$ & $\omega_{\curl F}^2$ \\
Divergence (\cref{thm:gauss_divergence}) &
  $\Omega\subseteq\mathbb{R}^n$ & $n$ & $\omega_F^{n-1}$ & $(\divg F)\,dx^1\!\wedge\!\cdots\!\wedge\! dx^n$ \\
\bottomrule
\end{tabular}
\end{center}
Each row's $\int_M d\omega = \int_{\partial M}\omega$ reads off as exactly that row's theorem.
The top row is \cref{sec:gauss_intuition}'s $n=1$ computation seen from the other side: there,
$f$ played the role of a $1$-dimensional vector field and the identity was reached via the
divergence theorem; here $f$ is instead a $0$-form and the same identity is the $m=1$ case of
\cref{thm:stokes_general} directly. These are two routes to the same fact, matching
\cref{ex:exterior_derivative_gradient_curl}'s remark that $df = \langle\nabla f,\cdot\rangle$ is
exactly the shape that makes this specialization work. The bottom row uses
$d\omega_F^{n-1} = (\divg F)\,dx^1\wedge\cdots\wedge dx^n$
(\cref{ex:exterior_derivative_gradient_curl}, item 3) together with the identification of
$\int_{\partial\Omega}\omega_F^{n-1}$ with the flux $\int_{\partial\Omega}\langle
F,\nu\rangle\,d\vol_{n-1}$ already noted in \cref{ch:pullbacks}.

The pattern is worth stating in words, because it is easy to lose under the notation: a
$k$-dimensional derivative measured \emph{inside} a region always equals the corresponding
$(k{-}1)$-dimensional quantity measured on its \emph{boundary}. Differentiation and taking a
boundary are, in this precise sense, the same operation viewed from two sides. That is also
exactly the content of $d^2=0$ (\cref{thm:exterior_derivative}\textbf{(3)}): a boundary has no
boundary, just as $d\omega$ has no further exterior derivative to take that isn't already zero.
\end{remark}

% Source: old_exams/examFS24.pdf (21 August 2024), Problem 3, parts 1--2, p. 10
% Extractor: Gemini 3.6 Flash (High)
% Correction: Claude Opus 5 (High) --- \operatorname{rot} -> \curl and J\Phi -> \Jac\Phi per the
% declared-macro rule; \iint -> \int, matching thm:green_theorem two pages above.
\begin{exercise}[Curl identity and Green's formula area preservation]
\label{ex:curl_identity_greens_formula}
Let $\Omega \subseteq \mathbb{R}^2$ be a bounded $C^1$ domain, and let $\Phi = (\Phi_1, \Phi_2) \in C^2(\overline{\Omega}, \mathbb{R}^2)$ be a vector-valued mapping, not assumed to be a diffeomorphism.
\begin{enumerate}[label=\textbf{(\alph*)}]
  \item Consider the vector field $V \in C^1(\overline{\Omega}, \mathbb{R}^2)$ defined by
  \[V := (\Phi_1 \partial_1 \Phi_2, \Phi_1 \partial_2 \Phi_2)\transp = \Phi_1 \nabla \Phi_2.\]
  Prove that $\curl V(x) = \det \Jac\Phi(x)$ for all $x \in \overline{\Omega}$.
  \item Let $\Psi \in C^2(\overline{\Omega}, \mathbb{R}^2)$ be another mapping such that $\Phi(x) = \Psi(x)$ for all $x \in \partial \Omega$. Prove that
  \[\int_\Omega \det \Jac\Phi(x) \, dx = \int_\Omega \det \Jac\Psi(x) \, dx.\]
\end{enumerate}
\exhint[Official hint]{Use Green's formula (i.e.\ 2D Stokes) on the plane.}
\exinfo{This exercise is taken from the exam of 21 August 2024, Problem 3. The original continues
with four further parts about $\Phi(\overline{B_1})$ when $\Phi$ fixes the unit circle.}
\end{exercise}

